% !TeX root = ../main.tex
\chapter{References}
\label{ch:references}

\section{Works the Pattern Names Are Taken From}

The architectural patterns named in Chapter~\ref{ch:patterns}, and the
advantages and disadvantages attributed to each, follow Sommerville. The
object-oriented patterns named in Chapter~\ref{ch:design-patterns} follow Gamma
and colleagues, commonly cited as the Gang of Four.

\begin{itemize}
	\item Ian Sommerville. \emph{Software Engineering}. Pearson. Chapter 6,
	      Architectural Design, covers \pattern{Pipe and Filter},
	      \pattern{Repository}, \pattern{Client-Server}, \pattern{Layered} and
	      \pattern{MVC}, and is the source for the properties of each that
	      Chapter~\ref{ch:patterns} tests this system against.
	\item Erich Gamma, Richard Helm, Ralph Johnson and John Vlissides.
	      \emph{Design Patterns: Elements of Reusable Object-Oriented
	      Software}. Addison-Wesley. The source for \textsc{Composite},
	      \textsc{Visitor}, \textsc{Adapter}, \textsc{Strategy},
	      \textsc{Command} and the five patterns considered and not adopted.
\end{itemize}

\section{Project Documents}

This document is the fourth of four and assumes the other three rather than
repeating them.

\begin{itemize}
	\item \emph{System Requirements Specification}, Team Pentateuch. Defines
	      every requirement identifier used here, along with the constraints
	      \textbf{C-n}, assumptions \textbf{A-n} and risks \textbf{R-n}. It is
	      the authority for the text of any requirement quoted in this
	      document.
	\item \emph{UML Design Document}, Team Pentateuch. Models the structure and
	      behaviour this document explains the reasoning behind. Its block
	      diagram is the decomposition that Chapter~\ref{ch:subsystems}
	      elaborates, and its class diagram chapter records the GENESYS schema
	      properties cited here.
	\item \emph{Tool Selection Document}, Team Pentateuch. Records which
	      external tools and interfaces the project builds on and why the
	      alternatives were set aside. Chapter~\ref{ch:patterns} treats those
	      choices as settled and reasons about the structure they imply.
\end{itemize}

\section{A Note on the GENESYS Schema}

The schema properties this document relies on, chiefly that a transition is an
entity rather than an attribute and that sequencing is carried by the fixed set
of behaviour constructs in \textbf{C-11}, were established by inspecting a
GENESYS project and the published interfaces of the Client API, not taken from
vendor documentation. They are recorded in the UML Design Document so that a
later reader can tell which parts of the design rest on observation rather than
on a specification.
